\begingroup\raggedright
\begin{thebibliography}{9}
\bibitem{BK} R. M. Bryant and L. G. Kov\'acs,
\emph{Tensor products of representations of finite groups},
Bull. London Math. Soc. \textbf{4} (1972), 133--135.
\href{https://archives.maths.anu.edu.au/people/Kovacs/K033.pdf}{Author archive}.
\bibitem{Ik} C. Ikenmeyer,
\emph{The Saxl conjecture and the dominance order},
Discrete Math. \textbf{338} (2015), 1970--1975.
\href{https://arxiv.org/abs/1410.6549v2}{arXiv:1410.6549v2}.
\bibitem{HR} N. Harman and C. Ryba,
\emph{A tensor-cube version of the Saxl conjecture},
Algebraic Combinatorics \textbf{6} (2023), 507--511.
\href{https://doi.org/10.5802/alco.267}{doi:10.5802/alco.267}.
\bibitem{BBS} C. Bessenrodt, C. Bowman, and L. Sutton,
\emph{Kronecker positivity and 2-modular representation theory},
Trans. Amer. Math. Soc. Ser. B \textbf{8} (2021), 1024--1055.
\href{https://doi.org/10.1090/btran/70}{doi:10.1090/btran/70}.
\bibitem{BB} C. Bessenrodt and C. Bowman,
\emph{Splitting Kronecker squares, 2-decomposition numbers, Catalan combinatorics, and the Saxl conjecture}, preprint, version 3 (2023).
\href{https://arxiv.org/abs/2202.03066v3}{arXiv:2202.03066v3}.
\bibitem{OAI} OpenAI,
\emph{A Cyclic Polytabloid Proof of Saxl's Conjecture}, September 24, 2026,
revision \texttt{fd4aeeb2ee4f}\allowbreak\texttt{c729c18d98444fed42fd0529eeeb}.
\href{https://github.com/openai/math/tree/fd4aeeb2ee4fc729c18d98444fed42fd0529eeeb/preprints/A-Cyclic-Polytabloid-Proof-of-Saxls-Conjecture-September-24-2026}{Pinned source}.
\bibitem{Universal} OpenAI,
\emph{Universal Tensor Squares for Symmetric Groups}, September 24, 2026,
same pinned revision as \cite{OAI}.
\href{https://github.com/openai/math/tree/fd4aeeb2ee4fc729c18d98444fed42fd0529eeeb/preprints/Universal-Tensor-Squares-for-Symmetric-Groups-September-24-2026}{Pinned source}.
\bibitem{I21} M. Kang,
\emph{Constructive signed dominance maps, their ranks, and wider-band limits},
project research note I21, October 9, 2026, Section~2.1.
Frozen SHA-256: \texttt{607c2c3296b6e020d8f0ae12526ee593}\allowbreak\texttt{ba3369e00fa5c4f393e53ee5b6936b6a}.
The construction used in that section is described in Section~\ref{sec:boundary} above.
\end{thebibliography}
\endgroup
